\documentclass[12pt]{article}

%============================================================
% IDIOMA Y CODIFICACIÓN
%============================================================
\usepackage[spanish, activeacute]{babel}
\usepackage[utf8]{inputenc}
\usepackage[T1]{fontenc}

%============================================================
% MATEMÁTICAS
%============================================================
\usepackage{amsmath,amssymb}

%============================================================
% GEOMETRÍA Y DISEÑO DE PÁGINA
%============================================================
\usepackage[margin=2.5cm]{geometry}
\setlength{\parindent}{0pt}

%============================================================
% LISTAS, ENLACES Y DIBUJOS
%============================================================
\usepackage{enumitem}
\usepackage{hyperref}
\hypersetup{colorlinks=true, linkcolor=black, urlcolor=blue, citecolor=blue}

%============================================================
% LISTADOS DE CÓDIGO
%============================================================
\usepackage{listings}
\usepackage[table]{xcolor}

\definecolor{codegreen}{rgb}{0,0.5,0}
\definecolor{codegray}{rgb}{0.5,0.5,0.5}
\definecolor{codepurple}{rgb}{0.5,0,0.5}
\definecolor{backcolour}{rgb}{0.95,0.95,0.95}

\lstdefinestyle{CodeStyle}{
    backgroundcolor=\color{backcolour},
    commentstyle=\color{codegreen},
    keywordstyle=\color{codepurple},
    numberstyle=\tiny\color{codegray},
    stringstyle=\color{codepurple},
    basicstyle=\ttfamily\footnotesize,
    breakatwhitespace=false,
    breaklines=true,
    captionpos=b,
    keepspaces=true,
    numbers=left,
    numbersep=5pt,
    showspaces=false,
    showstringspaces=false,
    showtabs=false,
    tabsize=2
}
\lstset{
    style=CodeStyle,
    literate=
        {á}{{\'a}}1 {é}{{\'e}}1 {í}{{\'i}}1 {ó}{{\'o}}1 {ú}{{\'u}}1
        {Á}{{\'A}}1 {É}{{\'E}}1 {Í}{{\'I}}1 {Ó}{{\'O}}1 {Ú}{{\'U}}1
        {ñ}{{\~n}}1 {Ñ}{{\~N}}1
        {ü}{{\"u}}1 {Ü}{{\"U}}1
}

%============================================================
% PSEUDOCÓDIGO
%============================================================
\usepackage[ruled,vlined]{algorithm2e}
\SetAlgorithmName{Algoritmo}{algoritmo}{Lista de algoritmos}

%============================================================
% ENCABEZADO (sin logo: este documento va solo, sin más archivos)
%============================================================
\usepackage{fancyhdr}
\pagestyle{fancy}
\fancyhf{}
\fancyhead[L]{\textbf{\small Estructura de Datos}}
\fancyhead[R]{\textbf{\small UNIVERSIDAD DE MAGALLANES}\\\textbf{\small Departamento de Ingeniería en Computación}}
\renewcommand{\headrulewidth}{0.4pt}
\setlength{\headheight}{30pt}
\fancyfoot[R]{\thepage}

\begin{document}

% --- Título ---
\begin{center}
    {\LARGE \textbf{¡Felicidades!}} \\[0.2cm]
    {\large Pudiste compilar tu primer documento en \LaTeX{}}
\end{center}
\vspace{0.3cm}

Si estás leyendo esto en un PDF generado por ti mismo, tu entorno de \LaTeX{} ya está listo para el curso. Aquí te dejo una pequeña muestra de por qué \LaTeX{} es genial.

\tableofcontents
\vspace{0.5cm}

%============================================================
\section{Fórmulas matemáticas}
%============================================================

Uno de los puntos fuertes de \LaTeX{} es lo natural que se ve la notación matemática. Por ejemplo, el teorema de Pitágoras:

\[
    a^2 + b^2 = c^2,
\]

o la fórmula general para resolver ecuaciones cuadráticas:

\[
    x = \frac{-b \pm \sqrt{b^2 - 4ac}}{2a}.
\]

%============================================================
\section{Notas al pie}
%============================================================

Como ayudante\footnote{Sí, las notas al pie son así de fáciles: solo se escribe \texttt{\textbackslash footnote\{...\}} justo después de la palabra que quieres aclarar.}, te recomiendo usarlas para comentarios que no quieres meter en medio del texto principal.

%============================================================
\section{Texto con estilo}
%============================================================

Puedes escribir en \textbf{negrita} cuando algo es importante, en \textit{cursiva} cuando quieres dar un énfasis más sutil, y en \texttt{fuente monoespaciada} cuando hablas de código, variables o comandos.

%============================================================
\section{Tablas}
%============================================================

\LaTeX{} también te permite armar tablas prolijas, con encabezados y filas destacadas con color:

\begin{table}[h]
    \centering
    \rowcolors{2}{white}{gray!15}
    \begin{tabular}{|l|l|l|l|}
        \hline
        \rowcolor{black}
        \textcolor{white}{\textbf{Fruta}} & \textcolor{white}{\textbf{Precio}} & \textcolor{white}{\textbf{Proveedor}} & \textcolor{white}{\textbf{Fecha de entrega}} \\
        \hline
        Manzana  & \$850  & Agrícola Los Andes  & 05/08/2026 \\
        Plátano  & \$600  & Frutas del Sur      & 08/08/2026 \\
        Naranja  & \$750  & Cítricos Patagonia  & 12/08/2026 \\
        Frutilla & \$1300 & Bayas Explora       & 15/08/2026 \\
        Kiwi     & \$980  & Huerto Austral      & 20/08/2026 \\
        \hline
    \end{tabular}
    \caption{Ejemplo de tabla con datos de proveedores de fruta}
\end{table}

%============================================================
\section{Bloques de código}
%============================================================

También puedes mostrar código fuente con numeración de líneas y resaltado de sintaxis:

\begin{lstlisting}[language=C, caption={Función de intercambio (swap)}, style=CodeStyle]
void swap(int *a, int *b) {
    int temp = *a;
    *a = *b;
    *b = temp;
}
\end{lstlisting}

%============================================================
\section{Pseudocódigo}
%============================================================

Y para cuando quieras describir un algoritmo sin atarte a un lenguaje en particular, aquí tienes uno que calcula el factorial de un número:

\begin{algorithm}[H]
    \caption{Cálculo del factorial de $n$}
    \KwIn{Un número entero $n \geq 0$}
    \KwOut{$n!$}
    $resultado \leftarrow 1$\;
    $i \leftarrow 1$\;
    \While{$i \leq n$}{
        $resultado \leftarrow resultado \times i$\;
        $i \leftarrow i + 1$\;
    }
    \Return{$resultado$}\;
\end{algorithm}

\end{document}
